\documentclass[../../main.tex]{subfiles}% 注意这里的文件路径不能用 ./main.tex ,否则用latexmk编译子文件会报错
\graphicspath{{\subfix{./image/}}} % 指定图片目录,后续可以直接使用图片文件名
% 注意这里的文件路径不能用 ../../image/ ,否则用latexmk编译子文件会报错

% 例如:
% \begin{figure}[H]
% \centering
% \includegraphics[scale=0.3]{图.png}
% \caption{}
% \label{figure:图}
% \end{figure}
% 注意:上述\label{}一定要放在\caption{}之后,否则引用图片序号会只会显示??.

\begin{document}

\section{Cantor(三分)集}

\begin{definition}
设$[0,1]\subset\mathbb{R}$，将$[0,1]$三等分，并移去中央三分开区间
\begin{align*}
I_{1,1}=\left(\frac{1}{3},\frac{2}{3}\right),
\end{align*}
记其留存部分为$F_1$，即
\begin{align*}
F_1=\left[0,\frac{1}{3}\right]\cup\left[\frac{2}{3},1\right]=F_{1,1}\cup F_{1,2};
\end{align*}
再将$F_1$中的区间$[0,1/3]$及$[2/3,1]$各三等分，并移去中央三分开区间
\begin{align*}
I_{2,1}=\left(\frac{1}{9},\frac{2}{9}\right) \quad \text{及} \quad I_{2,2}=\left(\frac{7}{9},\frac{8}{9}\right),
\end{align*}
记$F_1$中留存的部分为$F_2$，即
\begin{align*}
F_2&=\left[0,\frac{1}{9}\right]\cup\left[\frac{2}{9},\frac{1}{3}\right]\cup\left[\frac{2}{3},\frac{7}{9}\right]\cup\left[\frac{8}{9},1\right]\\
&=F_{2,1}\cup F_{2,2}\cup F_{2,3}\cup F_{2,4}.
\end{align*}
一般地说(归纳定义)，设所得剩余部分为$F_n$，则将$F_n$中每个（互不相交）区间三等分，并移去中央三分开区间，记其留存部分为$F_{n + 1}$，如此等等. 从而我们得到集合列$\{F_n\}$，其中
\begin{align*}
F_n=F_{n,1}\cup F_{n,2}\cup\cdots\cup F_{n,2^n} \quad (n = 1,2,\cdots).
\end{align*}
作点集$C = \bigcap_{n = 1}^{\infty}F_n$，我们称$C$为\textbf{Cantor(三分)集}. 
\end{definition}

\begin{theorem}[Cantor集的基本性质]\label{theorem:Cantor集的基本性质}
\begin{enumerate}[(1)]
\item $C$是非空有界闭集,因此是紧集.

\item $C = C'$,即$C$为完全集.

\item $C$无内点.

\item Cantor集的基数是$c$.

\item $[0,1]\setminus C$的长度的总和为1.
\end{enumerate}
\end{theorem}
\begin{proof}
\begin{enumerate}[(1)]
\item 因为每个$F_n$都是非空有界闭集，而且$F_n\supset F_{n + 1}$，所以根据Cantor闭集套定理，可知$C$不是空集（实际上，$F_n$ ($n = 1,2,\cdots$) 中每个闭区间的端点都是没有被移去的，即都是$C$中的点）. 显然，$C$是闭集.

\item 设$x\in C$，则$x\in F_n$ ($n = 1,2,\cdots$)，即对每个$n$，$x$属于长度为$1/3^n$的$2^n$个闭区间中的一个. 于是，对任一$\delta>0$，存在$n$，满足$1/3^n<\delta$，使得$F_n$中包含$x$的闭区间含于$(x - \delta, x + \delta)$. 此闭区间有两个端点，它们是$C$中的点且总有一个不是$x$. 这就说明$x$是$C$的极限点，故得$C'\supset C$. 由(i)知$C = C'$.

\item 设$x\in C$，给定任一区间$(x - \delta, x + \delta)$，取$2/3^n<\delta$. 因为$x\in F_n$，所以$F_n$中必有某个长度为$1/3^n$的闭区间$F_{n,k}$含于$(x - \delta, x + \delta)$. 然而，在构造$C$集的第$n + 1$步时，将移去$F_{n,k}$的中央三分开区间. 这说明$(x - \delta, x + \delta)$不含于$C$.

\item 事实上，将$[0,1]$中的实数按三进位小数展开，则Cantor集中点$x$与下述三进位小数集的元
\begin{align*}
x=\sum_{i = 1}^{\infty}\frac{a_i}{3^i}, \quad a_i = 0,2
\end{align*}
一一对应. 从而知$C$为连续基数集（与$(0,1]$的二进位小数比较）. 

\item 由$C$的定义可得
\begin{align*}
\sum_{n = 1}^{\infty}2^{n - 1}3^{-n}=1.
\end{align*}
\end{enumerate}

\end{proof}

\begin{definition}[类Cantor集]
设$\delta$是$(0,1)$内任意给定的数,考虑在$[0,1]$区间, 取$p=\frac{1+2\delta}{\delta}$，采用类似于Cantor集的构造过程:

第一步，移去长度为$\frac{1}{p}$的同心开区间；

第二步，在留存的两个闭区间的每一个中，又移去长度为$\frac{1}{p^2}$的同心开区间；

第三步，在留存的四个闭区间中再移去长度为$\frac{1}{p^3}$的同心区间. 继续此过程，可得一列移去的开区间，记其并集为$G$（开集）,
我们称$C_p=[0,1]\setminus G$为\textbf{类Cantor集}（当$p = 3$时，$C_p$就是Cantor（三分）集）.类Cantor集也称为\textbf{Harnack集}. 
\end{definition}
\begin{remark}
若要在$\mathbb{R}^n$的单位方体$[0,1]\times[0,1]\times\cdots\times[0,1]$中构造具有类似性质的集合，则只需取$C\times C\times\cdots\times C$（$C$是$[0,1]$中的类Cantor集）即可.
\end{remark}

\begin{theorem}[类Cantor集的基本性质]\label{theorem:类Cantor集的基本性质}
\begin{enumerate}[(1)]
\item $G$的总长度为$\delta$（$0<\delta<1$是任意给定的数）的稠密开集.

\item $C_p$是非空完全集，且没有内点.
\end{enumerate}
\end{theorem}
\begin{proof}
\begin{enumerate}[(1)]
\item 由$G$的定义可知$G$的总长度为
\begin{align*}
\sum_{n = 1}^{\infty}2^{n - 1}\left(\frac{1}{p}\right)^n=\frac{1}{p - 2}=\delta.
\end{align*}

\item 
\end{enumerate}

\end{proof}

\begin{definition}[Cantor函数]\label{definition:Cantor函数}
设 \(C\) 是 \([0,1]\) 中的 Cantor 集, 其中的点我们用三进位小数
\begin{align*}
x = 2\sum_{i = 1}^{\infty}\frac{\alpha_i}{3^i}, \quad \alpha_i = 0,1 \quad (i = 1,2,\cdots)
\end{align*}
来表示.

(i) 作定义在 \(C\) 上的函数 \(\varphi(x)\). 对于 \(x \in C\), 定义
\[
\varphi(x)=\varphi\left(2\sum_{i = 1}^{\infty}\frac{\alpha_i}{3^i}\right)=\sum_{i = 1}^{\infty}\frac{\alpha_i}{2^i}, \quad \alpha_i = 0,1 \quad (i = 1,2,\cdots).
\]

(ii) 作定义在 \([0,1]\) 上的 \(\varPhi(x)\). 对于 \(x \in [0,1]\), 定义
\[
\varPhi(x)=\sup\{\varphi(y):y\in C,y\leqslant x\}.
\]
我们称 \(\varPhi(x)\) 为 \textbf{Cantor函数}.
\end{definition}

\begin{theorem}[Cantor函数的性质]\label{theorem:Cantor函数的性质}
设$\varPhi(x)=\sup\{\varphi(y):y\in C,y\leqslant x\}$为Cantor函数, 则有下列性质:
\begin{enumerate}[(1)]
\item \(\varphi(C)=[0,1]\),即$\varphi$是满射.并且\(\varphi(x)\) 是 \(C\) 上的递增函数.

\item \(\varPhi(x)\) 是 \([0,1]\) 上的递增连续函数.此外, 在构造 Cantor 集的过程中所移去的每个中央三分开区间 \(I_{n,k}\) 上, \(\varPhi(x)\) 都是常数. 
\end{enumerate}
\end{theorem}
\begin{proof}
\begin{enumerate}[(1)]
\item 因为 \([0,1]\) 中的点可用二进位小数表示, 所以由$\varphi$的定义有 \(\varphi(C)=[0,1]\)..

下面证明$\varphi(x)$是$C$上的递增函数.设 \(\alpha_1,\alpha_2,\cdots,\beta_1,\beta_2,\cdots\) 是取 \(0\) 或 \(1\) 的数, 而且它们所表示的 \(C\) 中的数有下述关系:
\[
2\sum_{i = 1}^{\infty}\frac{\alpha_i}{3^i}<2\sum_{i = 1}^{\infty}\frac{\beta_i}{3^i}.
\]
若记 \(k = \min\{i:\alpha_i\neq\beta_i\}\), 则我们有
\begin{align*}
0&<\sum_{i = 1}^{\infty}\frac{\beta_i - \alpha_i}{3^i}=\frac{\beta_k - \alpha_k}{3^k}+\sum_{i > k}\frac{\beta_i - \alpha_i}{3^i}\\
&\leqslant\frac{\beta_k - \alpha_k}{3^k}+\sum_{i > k}\frac{2}{3^i}=\frac{\beta_k - \alpha_k + 1}{3^k}.
\end{align*}
由此可知 \((\alpha_k<\beta_k)\alpha_k = 0,\beta_k = 1\), 从而得到
\begin{align*}
\varphi\left(2\sum_{i = 1}^{\infty}\frac{\alpha_i}{3^i}\right)&=\sum_{i = 1}^{\infty}\frac{\alpha_i}{2^i}=\sum_{i = 1}^{k - 1}\frac{\alpha_i}{2^i}+\sum_{i = k}^{\infty}\frac{\alpha_i}{2^i}\\
&\leqslant\sum_{i = 1}^{k - 1}\frac{\beta_i}{2^i}+\sum_{i = k + 1}^{\infty}\frac{1}{2^i}=\sum_{i = 1}^{k - 1}\frac{\beta_i}{2^i}+\frac{1}{2^k}\\
&\leqslant\sum_{i = 1}^{k - 1}\frac{\beta_i}{2^i}+\sum_{i = k}^{\infty}\frac{\beta_i}{2^i}=\varphi\left(2\sum_{i = 1}^{\infty}\frac{\beta_i}{3^i}\right).
\end{align*}

\item 由(2)的结论及$\varPhi$的定义即得$\varPhi$的递增性.因为 \(\varPhi([0,1]) = [0,1]\), 所以由\refpro{Basis of Analytics-proposition:定义是区间的单调函数值域还是区间就必是连续函数}可知\(\varPhi(x)\) 是 \([0,1]\) 上的连续函数. 
\end{enumerate}

\end{proof}

\begin{example}
\(E\subset\mathbb{R}\) 是完全集当且仅当 \(E = \left(\bigcup_{n\geqslant 1}(a_n,b_n)\right)^c\)，其中 \((a_i,b_i)\) 与 \((a_j,b_j)\) \((i\neq j)\) 无公共端点。
\end{example}
\begin{proof}
{\heiti 必要性：}若 \(E\) 是完全集，则 \(E\) 是闭集。从而 \(E^c\) 是开集，它是 \(E^c\) 内构成区间的并集。这些构成区间相互之间是没有公共端点的，否则 \(E\) 中就会有孤立点了，这是不可能的。

{\heiti 充分性：}首先，由题设知 \(E\) 是闭集。其次，对任意的 \(x\in E\)，如果 \(x\notin E'\)，那么存在 \(\delta>0\)，使得 \((x - \delta,x + \delta)\cap E = \{x\}\)。这说明 \(x\) 是某两个开区间的端点，与假设矛盾。

\end{proof}

\begin{example}
设 \(E\subset\mathbb{R}^2\) 是完全集，则 \(E\) 是不可数集。
\end{example}
\begin{proof}
用反证法。假定 \(E = \{x_n\in\mathbb{R}^2: n = 1,2,\cdots\}\)。

(i) 选取 \(y_1\in E\setminus\{x_1\}\)，则点 \(x_1\) 到 \(y_1\) 的距离大于 \(0\)。存在以 \(y_1\) 为中心的闭正方形 \(Q_1\)，\(Q_1\cap E\) 是紧集。

(ii) 看 \(E\setminus\{x_2\}\)。因为 \(y_1\) 是 \(E\setminus\{x_2\}\) 的极限点，所以 \(\mathring{Q}_1\cap (E\setminus\{x_2\})\neq\varnothing\)。又取 \(y_2\in\mathring{Q}_1\cap (E\setminus\{x_2\})\)，并作以 \(y_2\) 为中心的闭正方形 \(Q_2\)：\(Q_2\subset Q_1\)，\(x_1\notin Q_2\)，\(x_2\notin Q_2\)，可知 \((Q_1\cap E)\supset (Q_2\cap E)\) 是紧集。如此继续做下去，可得有界闭集套列 \(\{Q_n\cap E\}\)：\((Q_{n - 1}\cap E)\supset (Q_n\cap E)\) \((n\in\mathbb{N})\)，而且 \(x_1,x_2,\cdots,x_n\) 不在其内。我们有
\begin{align*}
\bigcap_{n = 1}^{\infty}(Q_n\cap E)=\varnothing,
\end{align*}
导致矛盾。 

\end{proof}

\begin{proposition}
任一非空完全集的基数均为 \(c\)。
\end{proposition}
\begin{proof}
证明见那汤松著《实变函数论》的上册，有高等教育出版社出版的中译本，1955 年.

\end{proof}

\begin{example}
设 \(E = \left\{x\in[0,1]: x = \sum_{n = 1}^{\infty}a_n/10^n, a_n = 2 \text{ 或 } 7\right\}\)，我们有

(i) \(E\) 是闭集； 

(ii) \(\overline{E}=c\)； 

(iii) \(E\) 在 \([0,1]\) 中不稠密。
\end{example}
\begin{proof}
(i) 若有 \(\{x_m\}\subset E\)：\(x_m\rightarrow x (m\rightarrow\infty)\)，则
\begin{align*}
x = \sum_{n = 1}^{\infty}b_n/10^n \quad (b_n = 0,1,2,\cdots,9).
\end{align*}
如果 \(|x_m - x|<10^{-p}\)，那么在 \(x\in E\) 时，\(b_n = 2 \text{ 或 } 7 (n = 1,2,\cdots,p - 1)\)。这说明 \(E\) 是闭集。

(ii) 与 \(0\) 和 \(1\) 组成的数列类似，\(\overline{E}=c\)。

(iii) 注意到 \(E\cap(0.28,0.7)=\varnothing\)，故 \(E\) 不是稠密集。 

\end{proof}





\end{document}